% Part: first-order-logic
% Chapter: proof-systems
% Section: axiomatic-deduction

\documentclass[../../../include/open-logic-section]{subfiles}

\begin{document}

\iftag{FOL}
      {\olfileid{fol}{prf}{axd}}
      {\olfileid{pl}{prf}{axd}}

\olsection{स्वयंसिद्धकीय \usetoken{P}{derivation}}

स्वयंसिद्धकीय !!{derivation}s या सर्वांत जुन्या आणि सोप्या तार्किक
!!{derivation} पद्धती आहेत. त्यांतील !!{derivation}s म्हणजे केवळ
!!{sentence}s चे अनुक्रम असतात. !!{sentence}s चा अनुक्रम तेव्हा योग्य
!!{derivation} मानला जातो, जेव्हा त्यातील प्रत्येक !!{sentence}~$!A$ पुढील
अटींपैकी एक अट पूर्ण करते:
\begin{enumerate}
\item $!A$ हे स्वयंसिद्धक असते, किंवा
\item $!A$ हे दिलेल्या !!{sentence}s च्या संच~$\Gamma$ चा !!a{element} असते, किंवा
\item $!A$ ला निगमन नियमामुळे समर्थन मिळते.
\end{enumerate}
$!A$ हे स्वयंसिद्धक असण्यासाठी त्याचे रूप काही ठरावीक !!{sentence}
रूपबंधांपैकी एक असावे लागते. प्रथम-क्रम तर्कशास्त्रासाठी समाधानकारक
(निर्दोष आणि संपूर्ण) !!{derivation} पद्धत देणारे स्वयंसिद्धक-रूपबंधांचे
अनेक संच आहेत. त्यांपैकी काही संच, ते ज्या संयोजकांना नियंत्रित करतात
त्यांनुसार मांडलेले असतात; उदाहरणार्थ, पुढील रूपबंध
\[
!A \lif (!B \lif !A) \qquad !B \lif (!B \lor !C) \qquad (!B \land !C) \lif !B
\]
हे $\lif$, $\lor$ आणि~$\land$ यांना नियंत्रित करणारे सर्वसाधारण स्वयंसिद्धक
आहेत. काही स्वयंसिद्धकीय पद्धतींमध्ये स्वयंसिद्धकांची संख्या कमीत कमी
ठेवण्याचे उद्दिष्ट असते. कोणते संयोजक आदिम म्हणून स्वीकारले आहेत यानुसार,
एकाच स्वयंसिद्धकाची पद्धत मिळणेही शक्य असते.

निगमन नियम म्हणजे असे सोपाधिक विधान की जे !!a{derivation} मधील
!!a{sentence} समर्थित असण्यासाठी पुरेशी अट देते. विध्यात्मक अनुमान हा असा
एक अतिशय प्रचलित नियम आहे: $!A$ आणि $!A \lif !B$ यांना आधीच समर्थन मिळाले
असेल, तर $!B$ लाही समर्थन मिळते, असे तो सांगतो. म्हणजे, !!a{derivation}
मध्ये !!{sentence}~$!B$ असलेल्या ओळीला समर्थन मिळते, मात्र $!A$ आणि
$!A \lif !B$ ही दोन्ही सूत्रे (काही !!{sentence}~$!A$ साठी) त्या
!!{derivation} मध्ये~$!B$ च्या आधी आलेली असली पाहिजेत.

स्वयंसिद्धकीय !!{derivation}s वर आधारित $\Proves$ संबंधाची व्याख्या
पुढीलप्रमाणे केली जाते: $\Gamma \Proves !A$ तेव्हा आणि केवळ तेव्हाच, जेव्हा
अशी !!a{derivation} असते की तिच्या शेवटच्या सूत्राच्या जागी
!!{sentence}~$!A$ असते (आणि वरच्या~(2) मुळे समर्थन मिळालेल्या त्या
!!{derivation} मधील !!{sentence}s चा संच म्हणजे $\Gamma$ असे घेतले जाते).
~$!A$ ची अशी !!a{derivation} असेल की तिच्यासाठी~$\Gamma$ रिक्त असेल, म्हणजेच
त्या !!{derivation} मधील प्रत्येक !!{sentence} ला (1) किंवा~(3) मुळे समर्थन
मिळाले असेल, तर~$!A$ हे प्रमेय असते. उदाहरणार्थ, $\Proves
!A  \lif (!B \lif (!B \lor !A))$ हे दाखवणारी !!a{derivation} पुढे दिली आहे:
\begin{derivation}
  1. & $!B \lif (!B \lor !A)$ \\
  2. & $(!B \lif (!B \lor !A)) \lif (!A  \lif (!B \lif (!B \lor !A)))$\\
  3. & $!A  \lif (!B \lif (!B \lor !A))$
\end{derivation}
ओळ~1 वरील !!{sentence} हे $!A \lif (!A \lor !B)$ या स्वयंसिद्धकाच्या
रूपाचे आहे ($!A$ आणि $!B$ यांच्या भूमिका उलट आहेत). ओळ~2 वरील वाक्य
$!A \lif (!B \lif !A)$ या स्वयंसिद्धकाच्या रूपाचे आहे. त्यामुळे दोन्ही
ओळींना समर्थन मिळते. ओळ~3 ला विध्यात्मक अनुमानामुळे समर्थन मिळते: त्याचा
संक्षेप $!D$ असा केल्यास, ओळ~2 चे रूप $!C \lif !D$ असे आहे, आणि त्यातील
$!C$ म्हणजे $!B \lif (!B \lor !A)$, म्हणजेच ओळ~1 होय.

$\Gamma \Proves \lfalse$ असेल, तर संच $\Gamma$ विसंगत असतो. संपूर्ण
स्वयंसिद्धकीय पद्धत कोणत्याही~$!A$ साठी $\lfalse \lif !A$ हेही सिद्ध करेल;
त्यामुळे $\Gamma$ विसंगत असेल, तर कोणत्याही~$!A$ साठी $\Gamma \Proves !A$.

तर्कशास्त्रासाठी स्वयंसिद्धकीय !!{derivation}s च्या पद्धती सर्वप्रथम गोटलोप
फ्रेग यांनी त्यांच्या १८७९ मधील \emph{Begriffsschrift} या ग्रंथात दिल्या;
म्हणून हा ग्रंथ आधुनिक तर्कशास्त्रातील पहिला ग्रंथ मानला जातो. ॲल्फ्रेड
नॉर्थ व्हाइटहेड आणि बर्ट्रंड रसेल यांच्या \emph{Principia Mathematica} मध्ये,
तसेच १९२० च्या दशकात डाव्हीट हिल्बर्ट आणि त्यांच्या विद्यार्थ्यांकडून, या
पद्धतींना परिपूर्ण रूप मिळाले. म्हणून त्यांना अनेकदा ``फ्रेग पद्धती'' किंवा
``हिल्बर्ट पद्धती'' म्हणतात. एखाद्या तर्कशास्त्रासाठी स्वयंसिद्धकीय पद्धत
शोधणे बहुधा सोपे असल्यामुळे या पद्धती अतिशय बहुपयोगी आहेत. !!{derivation}s
ची रचना अत्यंत सोपी असून त्यांत फक्त एक किंवा दोन निगमन नियम असतात;
त्यामुळे \emph{त्यांच्याविषयीच्या} गोष्टी सिद्ध करणेही तुलनेने सोपे असते. तथापि,
प्रत्यक्ष व्यवहारात त्यांचा वापर फार कठीण असतो; म्हणजेच सिद्धता शोधणे आणि
लिहिणे कठीण असते.

\end{document}
